\documentclass[11pt]{article}
\usepackage[margin=1in]{geometry}
\usepackage{xcolor}
\usepackage{tcolorbox}
\usepackage{hyperref}
\usepackage{enumitem}
\setlist{noitemsep,topsep=0pt}
\usepackage{booktabs}
\usepackage{longtable}

% Define OSU orange and AI blue colors
\definecolor{osaorange}{RGB}{220,115,38}
\definecolor{aiblue}{RGB}{30,144,255}

% Configure hyperref colors
\hypersetup{
    colorlinks=true,
    linkcolor=blue,
    urlcolor=blue,
    citecolor=blue
}

% Remove default title
\title{}
\author{}
\date{}

\begin{document}

% Orange header box with black outline
\begin{tcolorbox}[colback=osaorange,colframe=black,boxrule=2pt,arc=0pt,left=6pt,right=6pt,top=10pt,bottom=10pt]
\centering
\textcolor{white}{\textbf{\large Oregon State University}}\\
\textcolor{white}{\textbf{\large College of Health}}\\
\textcolor{white}{\textbf{\large H524 Introduction to Biostatistics}}\\
\textcolor{white}{\textbf{\large Fall 2025}}
\end{tcolorbox}

\vspace{8pt}

\noindent\textbf{Credit Hours:} 04\\
\textbf{Room:} online\\
\textbf{Schedule:}

\vspace{4pt}

\begin{tabular}{@{}ll@{}}
Lecture: & Monday, Wednesday 12:00-1:20pm\\
Lab: & Wednesday: 8:00 am -- 9:50 am
\end{tabular}

\vspace{8pt}

\noindent\textbf{Course Instructor:} Dr. John Molitor\\
\textbf{Office Location:} 153 Milam\\
\textbf{Office Phone:} 541-737-3834\\
\textbf{E-Mail:} \href{mailto:john.molitor@oregonstate.edu}{john.molitor@oregonstate.edu}\\
\textbf{Office Hours:} Wednesday, 10:00 am -- 11:30 am

\vspace{8pt}

\noindent\textbf{Course Description:}

Quantitative analysis and interpretation of health data including probability distributions, estimation of associations, and significance tests such as chi-squared, one-way ANOVA, and linear regression.

\textbf{\textcolor{aiblue}{AI-Enhanced Component:}} This course integrates generative AI tools as learning partners for statistical analysis. Students will learn to leverage AI for code generation, concept explanation, and problem-solving while developing critical skills in output verification and responsible AI use.

\vspace{8pt}

\noindent\textbf{Academic Calendar}

All students are subject to the registration and refund deadlines as stated in the Academic Calendar: \url{https://registrar.oregonstate.edu/osu-academic-calendar}

\vspace{8pt}

\noindent\textbf{Instructor Communication}

Please contact me privately Canvas inbox for matters of a personal nature.

\vspace{8pt}

\noindent\textbf{Expectations for Time and Participation}

This course combines approximately 120 total hours of instruction, online activities, and assignments for 4 credits.

\vspace{8pt}

\noindent\textbf{Prerequisites}

H 524 assumes basic knowledge of statistical concepts and procedures equivalent to that of an undergraduate statistics course, and MTH 105 (Introduction to Contemporary Mathematics) or MTH 111 (College Algebra) or a higher mathematics course.

\vspace{8pt}

\noindent\textbf{Learning Resources}

Resources include, but are not limited to:
\begin{itemize}[leftmargin=20pt]
\item Lecture Notes
\item Reading Assignments
\item Discussions in class and on Canvas
\item Problem Sets
\item In-office help during office hours from both course instructor and lab instructors
\item \textcolor{aiblue}{\textbf{AI Tools:}} Access to ChatGPT, Claude, or GitHub Copilot (optional but recommended)
\item \textcolor{aiblue}{\textbf{AI Prompt Templates:}} Curated prompts for statistical analysis tasks
\end{itemize}

\vspace{8pt}

\noindent\textbf{Textbook}

Principles of Biostatistics, 2nd edition, 2022\\
Marcello Pagano and Kimberlee Gauvreau

\vspace{8pt}

\noindent\textbf{Measurable Student Learning Outcomes}

\begin{itemize}[leftmargin=20pt]
\item Select and generate graphical and numerical summaries of data.
\item Use principles of statistical inference to make conclusions about populations from samples.
\item Communicate statistical findings to others.
\item Use computer software to conduct statistical analysis.
\item \textcolor{aiblue}{\textbf{AI Integration:} Effectively prompt AI tools for statistical tasks and critically evaluate their outputs.}
\item \textcolor{aiblue}{\textbf{Verification Skills:} Validate AI-generated code and results through independent checking.}
\end{itemize}

\vspace{8pt}

\noindent\textbf{Evaluation of Student Performance}

\textbf{Grading:} Tests will be graded according to the following scale. It is possible, however, that adjustments will be made to this scale based on class performance.

\vspace{4pt}

\begin{center}
\begin{tabular}{cc}
\toprule
Range & Grade \\
\midrule
94--100\% & A \\
89--93 & A- \\
84--88 & B+ \\
79--83 & B \\
74--78 & B- \\
69--73 & C+ \\
64--68 & C \\
59--63 & C- \\
54--58 & D+ \\
49--53 & D \\
44--48 & D- \\
Less than 44 & F \\
\bottomrule
\end{tabular}
\end{center}

\vspace{8pt}

\noindent\textbf{Evaluation:} 30\% homework (includes AI-enhanced problem sets), 30\% midterm exam, 30\% final exam, 10\% class participation and AI reflection exercises.

\vspace{8pt}

\noindent\textbf{Homework}

Homework will be assigned regularly in class, approximately once a week. \textcolor{aiblue}{\textbf{AI Enhancement:} Each homework will include:}
\begin{itemize}[leftmargin=20pt]
\item \textcolor{aiblue}{Core statistical problems (70\% of grade)}
\item \textcolor{aiblue}{AI verification task - validate AI-generated solutions (20\% of grade)}
\item \textcolor{aiblue}{Brief AI usage documentation (10\% of grade)}
\end{itemize}

\vspace{8pt}

\noindent\textbf{AI Use Policy}

\begin{tcolorbox}[colback=aiblue!10,colframe=aiblue,boxrule=1pt,arc=2pt,left=6pt,right=6pt,top=6pt,bottom=6pt]
\textbf{Permitted AI Use:}
\begin{itemize}[leftmargin=10pt]
\item AI tools are encouraged for learning, code generation, and concept exploration
\item All AI-generated content must be verified and properly attributed
\item Students must demonstrate understanding independent of AI assistance
\end{itemize}

\textbf{Required Documentation:}
\begin{itemize}[leftmargin=10pt]
\item Cite AI tools used (name, version, date)
\item Include key prompts for significant AI contributions
\item Document verification steps taken
\end{itemize}

\textbf{Academic Integrity:}
\begin{itemize}[leftmargin=10pt]
\item Submitting unverified AI output as your own work is academic dishonesty
\item You are responsible for all errors in AI-generated content
\item Exams will test your understanding with and without AI assistance
\end{itemize}
\end{tcolorbox}

\vspace{8pt}

\noindent\textbf{Course Schedule}

The proposed schedule is subject to change by the instructor. What follows is a general outline of the proposed topics for discussion as well as the proposed dates. However, if we run out of time, get in a particularly lengthy discussion, or cover a topic more quickly than expected, there will likely be changes to the syllabus.

\vspace{8pt}

\begin{center}
\begin{tabular}{llll}
\toprule
\textbf{Week} & \textbf{Topic} & \textbf{Reading} & \textbf{\textcolor{aiblue}{AI Task}} \\
\midrule
Week 0 & Introduction to R & Pagano: & \textcolor{aiblue}{Configure AI tools;} \\
 & Data Presentation & Chapters 1, 2, 3 & \textcolor{aiblue}{Generate basic EDA code} \\
 & Numerical Summary Measures & & \\
\midrule
Week 1 & Introduction to R (cont.) & Pagano: & \textcolor{aiblue}{Compare AI plot} \\
 & Data Presentation & Chapters 1, 2, 3, 4 & \textcolor{aiblue}{suggestions;} \\
 & Numerical Summary Measures & & \textcolor{aiblue}{Document edits} \\
 & Rates and Standardization & & \\
\midrule
Week 2 & Probability & Pagano: Ch. 6 & \textcolor{aiblue}{Use AI for probability} \\
 & & (6.1--6.4.2, 6.5) & \textcolor{aiblue}{calculations; Verify with R} \\
\midrule
Week 3 & Theoretical Probability & Pagano: Ch. 7 & \textcolor{aiblue}{AI-scaffold} \\
 & Distributions & (7.1, 7.2, 7.4), 8 & \textcolor{aiblue}{CLT simulation} \\
 & Sampling Distribution of Mean & & \\
\midrule
Week 4 & Confidence Intervals & Pagano: Ch. 9 & \textcolor{aiblue}{Generate CI code;} \\
 & & & \textcolor{aiblue}{Compare interpretations} \\
\midrule
Week 5 & Hypothesis Testing & Pagano: Ch. 10 & \textcolor{aiblue}{Critique AI p-value} \\
 & Midterm exam & & \textcolor{aiblue}{explanations} \\
\midrule
Week 6 & Power & Pagano: & \textcolor{aiblue}{AI power analysis;} \\
 & Comparison of Two Means & Ch. 11, 12 & \textcolor{aiblue}{Validate results} \\
 & Analysis of Variance & & \\
\midrule
Week 7 & Nonparametric Methods & Pagano: Ch. 13 & \textcolor{aiblue}{When to use parametric} \\
 & & & \textcolor{aiblue}{vs non-parametric} \\
\midrule
Week 8 & Catch-up from Weeks 6 \& 7 & Pagano: & \textcolor{aiblue}{AI for OR/RR} \\
 & Contingency Tables & Ch. 15, 16 & \textcolor{aiblue}{interpretation} \\
 & Multiple Contingency Tables & & \\
\midrule
Week 9 & Correlation & Pagano: & \textcolor{aiblue}{Generate diagnostic} \\
 & Simple Linear Regression & Ch. 17, 18 & \textcolor{aiblue}{plots; Check assumptions} \\
\midrule
Week 10 & Simple Linear Regression & Pagano: Ch. 18 & \textcolor{aiblue}{Complete regression} \\
 & (continued) & & \textcolor{aiblue}{report with AI assistance} \\
\bottomrule
\end{tabular}
\end{center}

\textbf{Final Exam:} Available during finals week.

\vspace{8pt}

\noindent\textbf{Lab Sessions - AI Integration}

Wednesday labs will include:
\begin{itemize}[leftmargin=20pt]
\item Traditional R programming exercises
\item \textcolor{aiblue}{\textbf{AI-Enhanced Activities:}}
\begin{itemize}
\item Weeks 0-2: AI tool setup and basic prompting
\item Weeks 3-5: Using AI for statistical inference workflows
\item Weeks 6-7: AI-assisted method selection
\item Weeks 8-10: Building complete analyses with AI support
\end{itemize}
\end{itemize}

\vspace{8pt}

\noindent\textbf{Exam Format Updates}

\begin{itemize}[leftmargin=20pt]
\item \textbf{Midterm (Week 5):}
\begin{itemize}
\item Traditional closed-book section (50\%)
\item Open-resource section with AI permitted (50\%) - must document AI use
\end{itemize}
\item \textbf{Final Exam:}
\begin{itemize}
\item Comprehensive
\item Mix of AI-assisted and independent work
\item Includes AI output evaluation questions
\end{itemize}
\end{itemize}

\vspace{8pt}

\noindent\textbf{Diversity Statement}

The College of Public Health and Human Sciences strives to create an affirming climate for all students including underrepresented and marginalized individuals and groups. Diversity encompasses differences in age, color, ethnicity, national origin, gender, physical or mental ability, religion, socioeconomic background, veteran status, sexual orientation, and marginalized groups. We believe diversity is the synergy, connection, acceptance, and mutual learning fostered by the interaction of different human characteristics.

\vspace{8pt}

\noindent\textbf{Expectations for Student Conduct}

The Student Conduct Code establishes community standards and procedures necessary to maintain and protect an environment conducive to learning, in keeping with the educational objectives of Oregon State University. This code is based on the assumption that all persons must treat one another with dignity and respect in order for scholarship to thrive. For the full Student Conduct Code see \url{http://studentlife.oregonstate.edu/sites/studentlife.oregonstate.edu/files/code_of_student_conduct.pdf}

Academic or Scholarly Dishonesty is prohibited and considered a serious violation of the Student Conduct Code. It is defined as an act of deception in which a Student seeks to claim credit for the work or effort of another person, or uses unauthorized materials or fabricated information in any academic work or research, either through the Student's own efforts or the efforts of another. For specifics related to offenses proscribed by the University see: \url{http://oregonstate.edu/studentconduct/offenses-0}

\textcolor{aiblue}{\textbf{AI-Specific Academic Integrity:} Using AI tools without proper attribution, submitting unverified AI output as original work, or using AI to fabricate data or references constitutes academic dishonesty.}

\vspace{8pt}

\noindent\textbf{Religious Holiday Statement}

Oregon State University strives to respect all religious practices. If you have religious holidays that are in conflict with any of the requirements of this class, please see me immediately so that we can make alternative arrangements.

\vspace{8pt}

\noindent\textbf{Students with Documented Disabilities}

Accommodations for students with disabilities are determined and approved by Disability Access Services (DAS). If you, as a student, believe you are eligible for accommodations but have not obtained approval please contact DAS immediately at 541-737-4098 or at \url{http://ds.oregonstate.edu}. DAS notifies students and faculty members of approved academic accommodations and coordinates implementation of those accommodations. While not required, students and faculty members are encouraged to discuss details of the implementation of individual accommodations.

\end{document}